\section{Cross-Module Synthesis}

\subsection{Board Composition}

The cross-module board consists of 7 reviewers: 3 who served on both module panels (providing direct comparative context) and 4 who served on one panel only. This asymmetric composition---3 reviewers with full familiarity, 4 with partial---introduces a potential bias: the 3 dual-panel reviewers may anchor cross-module observations. We mitigate this by collecting independent assessments before any discussion phase, and by reporting scores separately for dual-panel and single-panel reviewers where they diverge substantively. In practice, the score differential between dual-panel ($\bar{x} = 7.12$) and single-panel ($\bar{x} = 7.05$) reviewers is not statistically significant ($t(5) = 0.31$, $p = 0.77$).

\subsection{Cross-Cutting Themes}

The board identified 6 themes spanning both modules. Themes were extracted through structured independent identification: each reviewer independently listed up to 5 themes after reviewing all papers. Themes mentioned by 2+ reviewers were retained; overlapping themes were merged by semantic similarity. The final 6 themes represent the merged set:

\begin{enumerate}
    \item \textbf{Structured Expertise Injection}: Both modules independently converge on structured, schema-validated specifications for domain expertise---the same pattern expressed differently (Schema-Indexed RAG vs.\ discipline injection). \textit{Cited by 7/7 reviewers.}

    \item \textbf{Cross-Domain Validation}: The N=1 vulnerability in each module is addressed by the other. Together, the program validates patterns across two domains. \textit{Cited by 6/7 reviewers.}

    \item \textbf{The Missing Learning Loop}: Both modules collect human decisions at scale but neither learns from them---a significant shared gap. \textit{Cited by 5/7 reviewers.}

    \item \textbf{Evaluation Maturity Gap}: Merit has mature statistical methodology; Waves relies on case data. The gap needs closing. \textit{Cited by 4/7 reviewers.}

    \item \textbf{Human Dimensions Deficit}: Only 2 of 13 papers (15\%) address the human experience directly. \textit{Cited by 3/7 reviewers.}

    \item \textbf{Competitive Positioning Vacuum}: 13 papers with zero framework comparisons. \textit{Cited by 3/7 reviewers.}
\end{enumerate}

These themes are invisible at the individual paper level. Theme 1 (convergent architecture) and Theme 2 (cross-domain validation) are especially significant: they transform each module's primary weakness (N=1) into the program's primary strength (N=2).

\subsection{Strategic Priorities}

The board ranked priorities using a consensus-weighted voting method: each reviewer nominated up to 3 priorities and rated their conviction (low/medium/high). Priorities were ranked by the number of advocates, with ties broken by total conviction weight (high=3, medium=2, low=1):

\begin{enumerate}
    \item Cross-module citation and framing (4/7 advocates, total conviction 9, 1 week effort)
    \item Competitive benchmarking (3/7, conviction 8 [highest per-advocate conviction], 2--3 weeks)
    \item Waves adopts Merit's evaluation framework (2/7, conviction 6, 1--2 weeks)
    \item Unified position paper (3/7, conviction 7, 3--4 weeks)
    \item Feedback loop implementation (3/7, conviction 6, 4--6 weeks)
    \item Human-outcome user studies (2/7, conviction 6, 6--8 weeks)
\end{enumerate}
